% ============================================================
%  يُترجَم بـ XeLaTeX  (في Overleaf: Menu → Compiler → XeLaTeX)
%  الخط: Amiri (متوفر في Overleaf، أو حمّله من github.com/aliftype/amiri)
% ============================================================
\documentclass[12pt,a4paper]{article}
\usepackage[a4paper,margin=1.7cm,top=2.3cm,bottom=2.9cm,headheight=15pt,footskip=1.5cm]{geometry}
\usepackage{amsmath,amssymb}
\usepackage{icomma}
\usepackage[english,bidi=default,layout=lists]{babel}
\babelprovide[import,main]{arabic}
\babelfont{rm}[Scale=1.08]{Amiri}
\babelfont{sf}[Scale=1.08]{Amiri}
\usepackage[most]{tcolorbox}
\usepackage{tikz}
\usetikzlibrary{arrows.meta,calc,shadings}
\usepackage{array}
\usepackage{colortbl}
\usepackage{fancyhdr}
\usepackage{titlesec}
\usepackage{enumitem}
\usepackage{collcell}
\usepackage{needspace}
\usepackage{etoolbox}
\newcolumntype{C}{>{\collectcell\ar}c<{\endcollectcell}}
\pretocmd{\section}{\Needspace{15\baselineskip}}{}{}
\pretocmd{\subsection}{\Needspace{11\baselineskip}}{}{}

\definecolor{cDef}{HTML}{1F4E9C}
\definecolor{cThm}{HTML}{B3261E}
\definecolor{cRem}{HTML}{2E7D32}
\definecolor{cEx}{HTML}{C76B00}
\definecolor{cMeth}{HTML}{6A1B9A}
\definecolor{cSol}{HTML}{00695C}
\definecolor{cSec}{HTML}{1F3A68}

\titleformat{\section}{\Large\bfseries\color{cSec}}{\thesection.}{0.6em}{}[{\color{cSec!60}\titlerule[1pt]}]
\titleformat{\subsection}{\large\bfseries\color{cThm!85!black}}{\thesubsection.}{0.5em}{}
\titlespacing*{\section}{0pt}{14pt}{8pt}
\titlespacing*{\subsection}{0pt}{10pt}{5pt}
\setlength{\parindent}{0pt}
\setlength{\parskip}{4pt}
\setlist{itemsep=2pt,topsep=3pt}

\newtcolorbox{arbox}[3][]{enhanced,
  colback=#2!5,colframe=#2!80!black,arc=2mm,boxrule=0.8pt,
  left=6pt,right=6pt,top=8pt,bottom=5pt,
  fonttitle=\bfseries,coltitle=white,
  title={\foreignlanguage{arabic}{#3}},
  attach boxed title to top right={xshift=-5mm,yshift=-3mm},
  boxed title style={colback=#2!80!black,arc=1.5mm,boxrule=0pt},
  before upper={\selectlanguage{arabic}},
  before skip=8pt,after skip=8pt,breakable,#1}
\newenvironment{BOX}[2]{\begin{otherlanguage}{english}\begin{arbox}{#1}{#2}}{\end{arbox}\end{otherlanguage}}
\newenvironment{tarif}[1][تعريف]{\begin{BOX}{cDef}{#1}}{\end{BOX}}
\newenvironment{khassia}[1][خاصية]{\begin{BOX}{cThm}{#1}}{\end{BOX}}
\newenvironment{molahada}[1][ملاحظة]{\begin{BOX}{cRem}{#1}}{\end{BOX}}
\newenvironment{mithal}[1][مثال]{\begin{BOX}{cEx}{#1}}{\end{BOX}}
\newenvironment{tariqa}[1][طريقة]{\begin{BOX}{cMeth}{#1}}{\end{BOX}}

\newcounter{exo}
\newenvironment{tamrin}[1]{\stepcounter{exo}\begin{BOX}{cSec}{التمرين \theexo: #1}}{\end{BOX}}
\newenvironment{hal}{\Needspace{12\baselineskip}\begin{otherlanguage}{english}\begin{arbox}[breakable]{cSol}{حل التمرين \theexo}}{\end{arbox}\end{otherlanguage}}

\newenvironment{LTR}{\begin{otherlanguage}{english}}{\end{otherlanguage}}
\newcommand{\ar}[1]{\foreignlanguage{arabic}{#1}}

\renewcommand{\dots}{\mathinner{\ldotp\ldotp\ldotp}}
\renewcommand{\ldots}{\mathinner{\ldotp\ldotp\ldotp}}
\renewcommand{\cdots}{\mathinner{\cdotp\cdotp\cdotp}}
\newcommand{\N}{\mathbb{N}}
\newcommand{\Z}{\mathbb{Z}}
\newcommand{\D}{\mathbb{D}}
\newcommand{\Q}{\mathbb{Q}}
\newcommand{\R}{\mathbb{R}}
\newcommand{\pgcd}{\operatorname{PGCD}}
\newcommand{\ppcm}{\operatorname{PPCM}}

% \NLseg{a}{b}{[}{]} ; \NLright{a}{[} ; \NLleft{b}{]}  (interval on a number line)
\newcommand{\NLseg}[4]{\begin{LTR}\begin{tikzpicture}[baseline=-0.6ex]
 \draw[-{Stealth[length=2mm]},gray!70](0,0)--(3.4,0);
 \draw[line width=2pt,cThm!80!black](0.7,0)--(2.7,0);
 \node[cThm!80!black,font=\LARGE] at (0.7,0){$#3$};
 \node[cThm!80!black,font=\LARGE] at (2.7,0){$#4$};
 \node[below=3pt,font=\footnotesize] at (0.7,0){$#1$};
 \node[below=3pt,font=\footnotesize] at (2.7,0){$#2$};
\end{tikzpicture}\end{LTR}}
\newcommand{\NLright}[2]{\begin{LTR}\begin{tikzpicture}[baseline=-0.6ex]
 \draw[-{Stealth[length=2mm]},gray!70](0,0)--(3.4,0);
 \draw[line width=2pt,cThm!80!black](1.3,0)--(3.3,0);
 \node[cThm!80!black,font=\LARGE] at (1.3,0){$#2$};
 \node[below=3pt,font=\footnotesize] at (1.3,0){$#1$};
\end{tikzpicture}\end{LTR}}
\newcommand{\NLleft}[2]{\begin{LTR}\begin{tikzpicture}[baseline=-0.6ex]
 \draw[-{Stealth[length=2mm]},gray!70](0,0)--(3.4,0);
 \draw[line width=2pt,cThm!80!black](0.05,0)--(2.1,0);
 \node[cThm!80!black,font=\LARGE] at (2.1,0){$#2$};
 \node[below=3pt,font=\footnotesize] at (2.1,0){$#1$};
\end{tikzpicture}\end{LTR}}

% ---------- social-media icons (pure TikZ) ----------
\definecolor{cFB}{HTML}{1877F2}
\definecolor{cTG}{HTML}{229ED9}
\definecolor{cIG1}{HTML}{FEDA75}
\definecolor{cIG2}{HTML}{D62976}
\definecolor{cIG3}{HTML}{4F5BD5}
\newcommand{\fbicon}{\begin{LTR}\begin{tikzpicture}[baseline=-0.75ex,x=1em,y=1em]
 \fill[cFB] (0,0) circle (0.55);
 \node[white,inner sep=0,font=\usefont{TU}{lmss}{bx}{n}\small] at (0.04,-0.02) {f};
\end{tikzpicture}\end{LTR}}
\newcommand{\igicon}{\begin{LTR}\begin{tikzpicture}[baseline=-0.75ex,x=1em,y=1em]
 \shade[lower left=cIG1,lower right=cIG2,upper left=cIG2,upper right=cIG3,rounded corners=0.28em] (-0.55,-0.55) rectangle (0.55,0.55);
 \draw[white,line width=0.09em,rounded corners=0.18em] (-0.36,-0.36) rectangle (0.36,0.36);
 \draw[white,line width=0.09em] (0,0) circle (0.17);
 \fill[white] (0.22,0.22) circle (0.05);
\end{tikzpicture}\end{LTR}}
\newcommand{\tgicon}{\begin{LTR}\begin{tikzpicture}[baseline=-0.75ex,x=1em,y=1em]
 \fill[cTG] (0,0) circle (0.55);
 \fill[white] (-0.33,0.0) -- (0.32,0.27) -- (0.2,-0.3) -- (0.0,-0.15) -- (-0.08,-0.28) -- (-0.08,-0.08) -- cycle;
 \fill[cTG!70!black] (-0.08,-0.08) -- (0.25,0.2) -- (0.0,-0.15) -- (-0.08,-0.28) -- cycle;
\end{tikzpicture}\end{LTR}}
% page number: black box, white two-digit number, at the right of the social row
\newcommand{\PageBox}{\begin{LTR}\tikz[baseline=(p.base)]\node[fill=black,text=white,inner sep=2.5pt,minimum width=1.6em,minimum height=1.3em,font=\usefont{TU}{lmss}{bx}{n}\small](p){\ifnum\value{page}<10 0\fi\arabic{page}};\end{LTR}}
% visual left-to-right order: Telegram | Instagram | Facebook | page box
\newcommand{\SocialRow}{\begin{otherlanguage}{english}%
 \ar{\textbf{الأستاذ سوايسية هشام للرياضيات}}\ \tgicon\hfill
 \textbf{prof\_hicham\_math}\ \igicon\hfill
 \ar{\textbf{الأستاذ سوايسية هشام للرياضيات}}\ \fbicon\hspace{2.2em}\PageBox\end{otherlanguage}}

\pagestyle{fancy}
\fancyhf{}
\renewcommand{\headrulewidth}{0.6pt}
\renewcommand{\footrulewidth}{0.4pt}
% \SetHeader{right head: document title}{left head: level -- stream}
\newcommand{\SetHeader}[2]{%
 \fancyhead[C]{\parbox{\textwidth}{\ar{#1\hfill #2}}}%
 \fancyfoot[C]{\parbox{\textwidth}{\small\SocialRow}}}
\usepackage{tkz-tab}
\usepackage{multicol}
% variation / sign table (LTR)
\newenvironment{TAB}{\par\begin{center}\begin{LTR}\begin{tikzpicture}[scale=0.95,every node/.style={font=\small}]}{\end{tikzpicture}\end{LTR}\end{center}}
% grid + axes, clipped: \begin{GR}{xmin}{xmax}{ymin}{ymax}{unit cm}
\newenvironment{GR}[5]{\par\begin{center}\begin{LTR}\begin{tikzpicture}[x=#5cm,y=#5cm,font=\footnotesize]
 \draw[gray!25,very thin] (#1,#3) grid (#2,#4);
 \draw[-{Stealth[length=2mm]}] (#1,0)--(#2+0.3,0) node[below left]{$x$};
 \draw[-{Stealth[length=2mm]}] (0,#3)--(0,#4+0.3) node[below left]{$y$};
 \node[below left] at (0,0){$O$};
 \begin{scope}\clip (#1,#3) rectangle (#2,#4);}
{\end{scope}\end{tikzpicture}\end{LTR}\end{center}}
\newcommand{\Ticks}[2]{\foreach \t in {#1} \node[below,font=\scriptsize] at (\t,0){$\t$}; \foreach \t in {#2} \node[left,font=\scriptsize] at (0,\t){$\t$};}
\newcommand{\bx}[1]{\ensuremath{\colorbox{cDef!12}{$\displaystyle #1$}}}
\newcommand{\TwoCol}[4]{\par\begin{LTR}\noindent\begin{minipage}[t]{#1\textwidth}\selectlanguage{arabic}#2\end{minipage}\hfill\begin{minipage}[t]{#3\textwidth}\selectlanguage{arabic}#4\end{minipage}\end{LTR}\par}
\newcommand{\RL}[4]{\TwoCol{#3}{#4}{#1}{#2}}  % \RL{w1}{right content}{w2}{left content}
\newcommand{\ok}{$\textcolor{cRem}{\checkmark}$}
\newcommand{\limn}{\displaystyle\lim_{n\to+\infty}}
% sub-questions, each on its own line
\newenvironment{sq}{\begin{itemize}[leftmargin=2.4em,itemsep=3pt]}{\end{itemize}}
\newcommand{\qa}{\item[(أ)]}\newcommand{\qb}{\item[(ب)]}\newcommand{\qc}{\item[(جـ)]}\newcommand{\qd}{\item[(د)]}\newcommand{\qe}{\item[(هـ)]}
\newcommand{\pa}{\par\smallskip\textbf{(أ)} }\newcommand{\pb}{\par\smallskip\textbf{(ب)} }\newcommand{\pc}{\par\smallskip\textbf{(جـ)} }\newcommand{\pd}{\par\smallskip\textbf{(د)} }\newcommand{\pe}{\par\smallskip\textbf{(هـ)} }
\newcommand{\QQ}[1]{\par\smallskip\textbf{#1.} }
\newenvironment{tmr}{\stepcounter{exo}\begin{BOX}{cSec}{تمرين محلول \theexo}}{\end{BOX}}
\newcommand{\ABC}[3]{\\[2pt]\begin{tabular}{p{0.28\textwidth}p{0.28\textwidth}p{0.28\textwidth}}(جـ) \ #3 & (ب) \ #2 & (أ) \ #1\end{tabular}}
\pretocmd{\section}{\Needspace{24\baselineskip}}{}{}
\newtcbox{\bxtbox}{on line,colback=cDef!12,colframe=cDef!12,boxrule=0pt,arc=1pt,boxsep=0pt,left=2pt,right=2pt,top=1pt,bottom=1pt}
\newcommand{\bxt}[1]{\foreignlanguage{english}{\bxtbox{\ar{#1}}}}
\usepackage{longtable}
\setlength{\LTpre}{4pt}\setlength{\LTpost}{0pt}
\newcommand{\Phase}[1]{\textbf{#1}}
\newcommand{\Fld}[2]{\textbf{#1:}\ #2}
% رأس المذكرة على نمط النموذج: إطار بعنوان «مذكرة رقم»، معلومات الحصة، ثم إطار الكفاءات
\newcommand{\ModhHead}[9]{%
\begin{LTR}
\begin{tcolorbox}[enhanced,colback=white,colframe=cThm!85!black,arc=4mm,boxrule=1pt,top=14pt,bottom=6pt,left=8pt,right=8pt,
 overlay={\node[fill=cThm!85!black,text=white,rounded corners=3mm,inner xsep=14pt,inner ysep=4pt,font=\large\bfseries] at ([yshift=0pt]frame.north){\ar{مذكرة رقم #1}};},
 before upper={\selectlanguage{arabic}}]
\begin{tabular}{@{}p{0.4\linewidth}p{0.55\linewidth}@{}}
\Fld{المستوى}{#2} & \Fld{الثانوية}{\dotfill}\\[3pt]
\Fld{المادة}{رياضيات} & \Fld{ميدان التعلم}{#3}\\[3pt]
\Fld{المدة}{#5} & \Fld{الوحدة التعلمية}{#4}\\[3pt]
\Fld{الأستاذ}{\dotfill} & \Fld{المحتوى المعرفي}{#6}\\
\end{tabular}
\end{tcolorbox}
\begin{tcolorbox}[enhanced,colback=white,colframe=cThm!85!black,arc=3mm,boxrule=0.8pt,sidebyside,sidebyside align=top,lefthand ratio=0.5,
 segmentation style={dashed,cThm!60},before upper={\selectlanguage{arabic}},before lower={\selectlanguage{arabic}}]
\textbf{الكفاءات المستهدفة:}\par #8
\tcblower
\textbf{الكفاءات القبلية:}\par #7
\end{tcolorbox}
\end{LTR}
\par\smallskip\textbf{الوسائل:} #9\par\smallskip}
\newenvironment{Sayr}{%
\begin{LTR}\renewcommand{\arraystretch}{1.25}
\begin{longtable}{|p{0.17\textwidth}|>{\centering\arraybackslash}p{0.06\textwidth}|p{0.53\textwidth}|>{\centering\arraybackslash}p{0.12\textwidth}|}
\hline\rowcolor{cThm!12}\ar{\textbf{توجيهات وتعاليق}} & \ar{\textbf{المدة}} & \ar{\textbf{سير الحصة}} & \ar{\textbf{مراحل الحصة}}\\\hline\endhead
\hline\endfoot}{\end{longtable}\end{LTR}}
% سطر في الجدول: \Row{مرحلة}{سير}{مدة}{توجيهات}
\newcommand{\Row}[4]{\ar{\small #4} & \ar{#3} & \ar{#2} & \cellcolor{cThm!8}\ar{\Phase{#1}}\\\hline}

\newcommand{\Lab}[2]{\foreignlanguage{english}{\tcbox[on line,colback=#1,colframe=#1,boxrule=0pt,arc=2pt,boxsep=1pt,left=5pt,right=5pt,top=1pt,bottom=1pt]{\ar{\textcolor{white}{\textbf{#2}}}}}}
\newcommand{\Num}[1]{\foreignlanguage{english}{\tikz[baseline=(c.base)]\node[fill=cDef,text=white,inner sep=2pt,font=\footnotesize\bfseries](c){#1};}}
\newcommand{\ck}{\ensuremath{{\color{cRem}\checkmark}}}
\newenvironment{Kad}[1]{\par\smallskip\Lab{#1}}{\par}
\SetHeader{\small\bfseries\color{cSec}مذكرات محور: المتتاليات العددية}{\small\color{cSec}السنة الثالثة ثانوي -- تسيير واقتصاد}
\begin{document}
\begin{LTR}
\begin{tcolorbox}[enhanced,colback=cThm!85!black,colframe=cThm!85!black,arc=3mm,halign=center,boxrule=0pt,top=12pt,bottom=12pt,before upper={\selectlanguage{arabic}}]
{\color{white}\fontsize{26}{32}\selectfont\bfseries مذكرات محور: المتتاليات العددية}\\[8pt]
{\color{white}\large السنة الثالثة ثانوي -- تسيير واقتصاد}\\[4pt]
{\color{white}\large 8 مذكرات (16 ساعة) -- السنة الدراسية 2026-2027}
\end{tcolorbox}
\end{LTR}
\begin{BOX}{cThm}{يتضمن هذا العمل}
\begin{enumerate}[label=\textbf{مذكرة \arabic{enumi}:},leftmargin=6em]
\item عموميات حول المتتاليات العددية.
\item المتتاليات الحسابية.
\item المتتاليات الهندسية.
\item الاستدلال بالتراجع.
\item المتتاليات المحدودة والرتيبة.
\item نهاية متتالية ، التقارب.
\item المتتاليات من الشكل $u_{n+1}=au_n+b$.
\item حل مشكلات اقتصادية بالمتتاليات.
\end{enumerate}
مراجع النشاطات: الكتاب المدرسي للسنة الثالثة ثانوي تسيير واقتصاد (باب المتتاليات ص 6--7).
\end{BOX}
\newpage
\ModhHead{01}{3 ت ا}{تحليل}{المتتاليات العددية}{02 سا}{عموميات حول المتتاليات العددية}%
{\begin{itemize}\item المتتاليات (2 ت ا) ، الدوال.\end{itemize}}%
{\begin{itemize}\item توليد متتالية بدستور صريح أو بعلاقة تراجعية ، وتمثيلها بيانياً.\end{itemize}}%
{المنهاج ، الكتاب المدرسي ، مراجع ، الحاسبة.}
\begin{Sayr}
\Row{مرحلة الانطلاق}{%
\textbf{التهيئة النفسية:} حساب توفير فيه $1000$ DA ، يزيد كل شهر بـ $2\,\%$: كيف نعبر عن الرصيد شهراً بعد شهر؟}{10 د}{وضعية اقتصادية.}
\Row{التشخيص والاكتشاف}{%
\Lab{cEx}{وضعية مقترحة}\par\ck\ $u_0=1000$ ، $u_{n+1}=1{,}02\,u_n$: كل رصيد يحسب من الذي قبله (علاقة تراجعية).\par
\ck\ $u_1=1020$ ، $u_2=1040{,}4$.\par
}{35 د}{الحاسبة.}
\Row{بناء المعارف}{%
\Lab{cDef}{تعريف}\ المتتالية العددية دالة من $\mathbb N$ نحو $\mathbb R$ ؛ $u_n$ حدها العام.\par\Lab{cMeth}{طرق التوليد}\par\ck\ صريحة: $u_n=n^2-4n$: $u_0=0$ ، $u_1=-3$ ، $u_2=-4$.\par
\ck\ تراجعية: $u_0=2$ ، $u_{n+1}=2u_n-1$: $u_1=3$ ، $u_2=5$ ، $u_3=9$.\par
}{50 د}{تمثيل النقط.}
\Row{}{%
\Lab{cRem}{ملاحظات}\ التمثيل البياني: نقط منعزلة $(n;u_n)$.\par\Lab{cThm!80!black}{حذار}\ في العلاقة التراجعية نحتاج الحد السابق.}{}{}
\Row{التقويم}{%
\Lab{cEx}{تطبيق}\par احسب $v_0$ ، $v_1$ ، $v_2$ حيث $v_n=\dfrac{n+1}{n+2}$.\par\textbf{الحل:} $\frac12$ ، $\frac23$ ، $\frac34$.\par\textbf{واجب منزلي:} تمارين من باب المتتاليات.\par\smallskip\textbf{ملاحظات حول سير الحصة:} \dotfill}{25 د}{تقويم فردي.}
\end{Sayr}
\newpage
\ModhHead{02}{3 ت ا}{تحليل}{المتتاليات العددية}{02 سا}{المتتاليات الحسابية}%
{\begin{itemize}\item عموميات المتتاليات.\end{itemize}}%
{\begin{itemize}\item التعرف على متتالية حسابية ، حساب حدها العام ومجموع حدود متتابعة.\end{itemize}}%
{المنهاج ، الكتاب المدرسي ، مراجع ، الحاسبة.}
\begin{Sayr}
\Row{مرحلة الانطلاق}{%
\textbf{التهيئة النفسية:} إيجار $30000$ DA سنوياً يزيد بـ $1500$ DA كل سنة: كم يصبح بعد $10$ سنوات؟}{10 د}{وضعية الإيجار.}
\Row{التشخيص والاكتشاف}{%
\Lab{cEx}{مناقشة نشاط 01 صفحة 6}\par\ck\ الفرق بين حدين متتابعين ثابت: متتالية حسابية.\par
}{35 د}{نشاط الكتاب.}
\Row{بناء المعارف}{%
\Lab{cDef}{تعريف}\ $u_{n+1}=u_n+r$ ($r$ الأساس).\par\Lab{cThm}{خاصية}\ $u_n=u_0+nr=u_p+(n-p)r$ ؛ المجموع $S=N\times\dfrac{a+b}2$ ($N$ عدد الحدود ، $a$ الأول ، $b$ الأخير).\par}{50 د}{الدستوران.}
\Row{}{%
\Lab{cEx}{مثال}\ $u_n=30000+1500n$: $u_{10}=45000$ DA ؛ مجموع إيجارات السنوات $u_0$ إلى $u_9$: $10\times\dfrac{30000+43500}2=367500$ DA.\par\Lab{cThm!80!black}{حذار}\ نعدّ الحدود بدقة.}{}{}
\Row{التقويم}{%
\Lab{cEx}{تطبيق}\par متتالية حسابية: $u_5=17$ و $u_{12}=38$. عيّن $r$ و $u_0$ ، ثم $u_0+\cdots+u_{12}$.\par\textbf{الحل:} $7r=21$ ، $r=3$ ، $u_0=2$ ؛ المجموع $13\times\dfrac{2+38}2=260$.\par\textbf{واجب منزلي:} تمارين من باب المتتاليات.\par\smallskip\textbf{ملاحظات حول سير الحصة:} \dotfill}{25 د}{تقويم فردي.}
\end{Sayr}
\newpage
\ModhHead{03}{3 ت ا}{تحليل}{المتتاليات العددية}{02 سا}{المتتاليات الهندسية}%
{\begin{itemize}\item المتتاليات الحسابية ، النسب المئوية.\end{itemize}}%
{\begin{itemize}\item التعرف على متتالية هندسية ، حساب حدها العام ومجموع حدود متتابعة ، ونمذجة الفوائد المركبة.\end{itemize}}%
{المنهاج ، الكتاب المدرسي ، مراجع ، الحاسبة.}
\begin{Sayr}
\Row{مرحلة الانطلاق}{%
\textbf{التهيئة النفسية:} رأس مال $100000$ DA بفائدة مركبة $5\,\%$ سنوياً: كم يصبح بعد $3$ سنوات؟}{10 د}{وضعية بنكية.}
\Row{التشخيص والاكتشاف}{%
\Lab{cEx}{وضعية مقترحة}\par\ck\ كل سنة نضرب في $1{,}05$: $C_{n+1}=1{,}05\,C_n$.\par
}{35 د}{مقارنة بالحسابية.}
\Row{بناء المعارف}{%
\Lab{cDef}{تعريف}\ $v_{n+1}=q\,v_n$ ($q$ الأساس).\par\Lab{cThm}{خاصية}\ $v_n=v_0q^n$ ؛ المجموع ($q\neq1$): $S=a\times\dfrac{1-q^N}{1-q}$ ($a$ الحد الأول ، $N$ عدد الحدود).\par}{50 د}{الحاسبة.}
\Row{}{%
\Lab{cEx}{مثال}\ $C_n=100000\times1{,}05^n$: $C_3=115762{,}5$ DA.\par\Lab{cRem}{ملاحظات}\ نسبة تطور $t\,\%$ تعني الضرب في $1+\frac t{100}$.\par}{}{}
\Row{}{%
\Lab{cThm!80!black}{حذار}\ الفائدة المركبة هندسية ، والبسيطة حسابية.}{}{}
\Row{التقويم}{%
\Lab{cEx}{تطبيق}\par متتالية هندسية: $v_1=6$ و $v_4=48$. عيّن $q$ و $v_0$ ، ثم $v_0+\cdots+v_5$.\par\textbf{الحل:} $q^3=8$ ، $q=2$ ، $v_0=3$ ؛ $S=3(2^6-1)=189$.\par\textbf{واجب منزلي:} تمارين من باب المتتاليات.\par\smallskip\textbf{ملاحظات حول سير الحصة:} \dotfill}{25 د}{تقويم فردي.}
\end{Sayr}
\newpage
\ModhHead{04}{3 ت ا}{تحليل}{المتتاليات العددية}{02 سا}{الاستدلال بالتراجع}%
{\begin{itemize}\item المتتاليات.\end{itemize}}%
{\begin{itemize}\item إثبات خاصية متعلقة بعدد طبيعي بالاستدلال بالتراجع.\end{itemize}}%
{المنهاج ، الكتاب المدرسي ، مراجع ، الحاسبة.}
\begin{Sayr}
\Row{مرحلة الانطلاق}{%
\textbf{التهيئة النفسية:} صف قطع دومينو: ماذا يكفي لتسقط كلها؟}{10 د}{مثال الدومينو.}
\Row{التشخيص والاكتشاف}{%
\Lab{cEx}{وضعية مقترحة}\ $u_0=2$ ، $u_{n+1}=2u_n+1$: $u_1=5$ ، $u_2=11$ ، $u_3=23$.\par\ck\ تخمين: $u_n=3\times2^n-1$.\par
}{35 د}{تخمين ثم برهان.}
\Row{بناء المعارف}{%
\Lab{cThm}{خاصية}\ لإثبات $P(n)$ من أجل كل $n\geqslant n_0$: \textbf{التحقق} من $P(n_0)$ ، ثم \textbf{الوراثة}: نفرض $P(n)$ ونبرهن $P(n+1)$.\par\Lab{cEx}{مثال}\ التحقق: $3\times1-1=2=u_0$ ؛ الوراثة: $u_{n+1}=2(3\times2^n-1)+1=3\times2^{n+1}-1$.\par}{50 د}{تحرير المراحل.}
\Row{}{%
\Lab{cThm!80!black}{حذار}\ لا ننسى مرحلة التحقق.}{}{}
\Row{التقويم}{%
\Lab{cEx}{تطبيق}\par برهن أنه من أجل كل $n\in\mathbb N$: $2^n\geqslant n+1$.\par\textbf{الحل:} $n=0$: $1\geqslant1$ ؛ $2^{n+1}=2\times2^n\geqslant2n+2\geqslant n+2$.\par\textbf{واجب منزلي:} تمارين من باب المتتاليات.\par\smallskip\textbf{ملاحظات حول سير الحصة:} \dotfill}{25 د}{تقويم فردي.}
\end{Sayr}
\newpage
\ModhHead{05}{3 ت ا}{تحليل}{المتتاليات العددية}{02 سا}{المتتاليات المحدودة والرتيبة}%
{\begin{itemize}\item الاستدلال بالتراجع.\end{itemize}}%
{\begin{itemize}\item إثبات أن متتالية محدودة ، ودراسة رتابتها.\end{itemize}}%
{المنهاج ، الكتاب المدرسي ، مراجع ، الحاسبة.}
\begin{Sayr}
\Row{مرحلة الانطلاق}{%
\textbf{التهيئة النفسية:} $v_n=\dfrac{2n+1}{n+1}$: هل تتجاوز الحدود العدد $2$ ؟}{10 د}{سؤال.}
\Row{التشخيص والاكتشاف}{%
\Lab{cEx}{وضعية مقترحة}\par\ck\ $v_n=2-\dfrac1{n+1}$: $1\leqslant v_n<2$ ، وتتزايد.\par
}{35 د}{كتابة جديدة.}
\Row{بناء المعارف}{%
\Lab{cDef}{تعريف}\ محدودة من الأعلى: $u_n\leqslant M$ ؛ من الأسفل: $u_n\geqslant m$ ؛ محدودة: الاثنان معاً.\par\Lab{cMeth}{الرتابة}\ إشارة $u_{n+1}-u_n$ ، أو $\dfrac{u_{n+1}}{u_n}$ مقارنة بـ $1$ (حدود موجبة) ، أو بالتراجع.\par}{50 د}{جدول الطرق.}
\Row{}{%
\Lab{cEx}{مثال}\ $u_0=0$ ، $u_{n+1}=\frac12u_n+1$: بالتراجع $0\leqslant u_n<2$ ؛ $u_{n+1}-u_n=1-\frac12u_n>0$: متزايدة.\par\Lab{cThm!80!black}{حذار}\ نبرهن الحصر قبل استعماله في الرتابة.}{}{}
\Row{التقويم}{%
\Lab{cEx}{تطبيق}\par ادرس رتابة $w_n=\dfrac{3^n}{n+1}$ ($n\geqslant1$).\par\textbf{الحل:} $\dfrac{w_{n+1}}{w_n}=\dfrac{3(n+1)}{n+2}>1$: متزايدة.\par\textbf{واجب منزلي:} تمارين 15 ، 17 ، 18 ، 19 من باب المتتاليات.\par\smallskip\textbf{ملاحظات حول سير الحصة:} \dotfill}{25 د}{تقويم فردي.}
\end{Sayr}
\newpage
\ModhHead{06}{3 ت ا}{تحليل}{المتتاليات العددية}{02 سا}{نهاية متتالية ، التقارب}%
{\begin{itemize}\item المتتاليات الرتيبة والمحدودة ، النهايات.\end{itemize}}%
{\begin{itemize}\item حساب نهاية متتالية ، ودراسة تقاربها.\end{itemize}}%
{المنهاج ، الكتاب المدرسي ، مراجع ، الحاسبة.}
\begin{Sayr}
\Row{مرحلة الانطلاق}{%
\textbf{التهيئة النفسية:} رصيد يزيد بـ $5\,\%$ سنوياً: هل يصبح «لا نهائياً»؟ وكمية تنقص بـ $20\,\%$ سنوياً؟}{10 د}{وضعيتان.}
\Row{التشخيص والاكتشاف}{%
\Lab{cEx}{وضعية مقترحة}\par\ck\ $1{,}05^n\to+\infty$ ، $0{,}8^n\to0$.\par
}{35 د}{الحاسبة.}
\Row{بناء المعارف}{%
\Lab{cThm}{خاصية}\par\ck\ $q^n$: $+\infty$ إذا $q>1$ ، $0$ إذا $-1<q<1$.\par
\ck\ كل متتالية متزايدة ومحدودة من الأعلى متقاربة (وكذلك متناقصة ومحدودة من الأسفل).\par
\ck\ إذا كان $u_n=f(n)$ و $\lim\limits_{x\to+\infty}f(x)=l$ فإن $\lim u_n=l$.\par
\Lab{cEx}{مثال}\ $\lim\dfrac{3n+1}{n+2}=3$ ؛ $\lim\left(5-2\times0{,}5^n\right)=5$.\par}{50 د}{النهايات المرجعية.}
\Row{}{%
\Lab{cThm!80!black}{حذار}\ متتالية محدودة ليست بالضرورة متقاربة.}{}{}
\Row{التقويم}{%
\Lab{cEx}{تطبيق}\par احسب نهاية $u_n=1000\times0{,}9^n+200$ وفسّر.\par\textbf{الحل:} $0{,}9^n\to0$: النهاية $200$.\par\textbf{واجب منزلي:} تمارين 16 ، 18 ، 27 ، 28 من باب المتتاليات.\par\smallskip\textbf{ملاحظات حول سير الحصة:} \dotfill}{25 د}{تقويم فردي.}
\end{Sayr}
\newpage
\ModhHead{07}{3 ت ا}{تحليل}{المتتاليات العددية}{02 سا}{المتتاليات من الشكل $u_{n+1}=au_n+b$}%
{\begin{itemize}\item المتتاليات الحسابية والهندسية ، النهايات.\end{itemize}}%
{\begin{itemize}\item دراسة متتالية $u_{n+1}=au_n+b$ بمتتالية مساعدة هندسية.\end{itemize}}%
{المنهاج ، الكتاب المدرسي ، مراجع ، الحاسبة.}
\begin{Sayr}
\Row{مرحلة الانطلاق}{%
\textbf{التهيئة النفسية:} ما الذي يحدث لكمية تنقص بالنصف ونضيف لها $3$ كل مرة؟}{10 د}{سؤال.}
\Row{التشخيص والاكتشاف}{%
\Lab{cEx}{مناقشة نشاط 04 صفحة 7}\ $u_0=1$ ، $u_{n+1}=0{,}5u_n+3$.\par\ck\ النقطة الصامدة: $l=0{,}5l+3$ ، $l=6$.\par
}{35 د}{نشاط الكتاب.}
\Row{بناء المعارف}{%
\Lab{cMeth}{طريقة}\ نضع $v_n=u_n-l$: $v_{n+1}=a\,v_n$ ، هندسية أساسها $a$ ، ثم $u_n=l+v_0a^n$.\par\Lab{cEx}{مثال}\ $v_n=u_n-6$ ، $v_0=-5$: $u_n=6-5\times0{,}5^n$ ، $\lim u_n=6$.\par}{50 د}{مخطط الطريقة.}
\Row{}{%
\Lab{cThm!80!black}{حذار}\ نبرهن أن $(v_n)$ هندسية بحساب $v_{n+1}$.}{}{}
\Row{التقويم}{%
\Lab{cEx}{تطبيق}\par $u_0=3$ ، $u_{n+1}=3u_n-4$. احسب $u_n$ ونهايتها.\par\textbf{الحل:} $l=2$ ، $v_n=u_n-2$ هندسية أساسها $3$ ، $v_0=1$: $u_n=2+3^n\to+\infty$.\par\textbf{واجب منزلي:} تمارين 15 ، 16 ، 22 من باب المتتاليات.\par\smallskip\textbf{ملاحظات حول سير الحصة:} \dotfill}{25 د}{تقويم فردي.}
\end{Sayr}
\newpage
\ModhHead{08}{3 ت ا}{تحليل}{المتتاليات العددية}{02 سا}{حل مشكلات اقتصادية بالمتتاليات}%
{\begin{itemize}\item المتتاليات $u_{n+1}=au_n+b$.\end{itemize}}%
{\begin{itemize}\item نمذجة وضعيات اقتصادية (إيداع ، سكان) بمتتالية وحلها.\end{itemize}}%
{المنهاج ، الكتاب المدرسي ، مراجع ، الحاسبة.}
\begin{Sayr}
\Row{مرحلة الانطلاق}{%
\textbf{التهيئة النفسية:} نودع $10000$ DA بفائدة $4\,\%$ سنوياً ونضيف $2000$ DA كل سنة: كم الرصيد بعد $5$ سنوات؟}{10 د}{وضعية بنكية.}
\Row{التشخيص والاكتشاف}{%
\Lab{cEx}{وضعية مقترحة}\par\ck\ $S_0=10000$ ، $S_{n+1}=1{,}04S_n+2000$.\par
}{35 د}{عمل بالأفواج.}
\Row{بناء المعارف}{%
\Lab{cEx}{الحل}\par\ck\ $l=1{,}04l+2000$ ، $l=-50000$ ؛ $v_n=S_n+50000$ هندسية أساسها $1{,}04$ ، $v_0=60000$.\par
\ck\ $S_n=60000\times1{,}04^n-50000$ ؛ $S_5\approx22999{,}17$ DA.\par
\Lab{cMeth}{طريقة}\ ترييض ، متتالية مساعدة ، حد عام ، حساب ، تفسير في سياق المشكلة.\par}{50 د}{تفسير النتائج.}
\Row{}{%
\Lab{cThm!80!black}{حذار}\ نحترم الوحدات ونقرب حسب المطلوب.}{}{}
\Row{التقويم}{%
\Lab{cEx}{تطبيق}\par مدينة سكانها $20000$ ، يغادرها $2\,\%$ سنوياً ويصلها $300$ ساكن: $P_{n+1}=0{,}98P_n+300$. عيّن $P_n$ ونهايتها.\par\textbf{الحل:} $l=15000$ ، $P_n=15000+5000\times0{,}98^n\to15000$.\par\textbf{واجب منزلي:} مسائل من باب المتتاليات (بكالوريا).\par\smallskip\textbf{ملاحظات حول سير الحصة:} \dotfill}{25 د}{نهاية المحور: فرض أو مراجعة.}
\end{Sayr}
\end{document}
